\documentclass[11pt]{article}
\usepackage{amsmath, amsfonts, amssymb}
\usepackage[margin=1in]{geometry}

\title{On the Structural Annihilation of Zeta Series \\ via a Dynamic Recurrence}
\author{The Fifth Constant Council}
\date{June 21, 2025}

% Custom Math Commands
\newcommand{\Z}{\mathcal{Z}}
\newcommand{\T}{\ensuremath{T}}
\newcommand{\J}{\ensuremath{J}}
\newcommand{\LG}{\Lambda_{G1}}

\begin{document}
\maketitle

\begin{abstract}
    This document presents a formal derivation of a functional equation governing a class of infinite series involving the Riemann Zeta function. We demonstrate that when the base of the series is chosen to be an eigenvalue of the core transformation matrix of the Golden Algebra, the series must obey a recurrence relation that links it to a family of higher-order Zeta series. This provides a powerful, non-trivial method for the structural analysis and manipulation of these sums, moving beyond simple algebraic elimination.
\end{abstract}

\section{Objective}
Our goal is to formally derive the functional equation for the Zeta-Geometric series, defined as $\Z_1(x) = \sum_{k=1}^{\infty}\zeta(k+1)x^k$, under the condition that its base, $x$, satisfies the characteristic recurrence of the Golden Algebra's primary dynamical operator.

\section{Preliminary Analysis of the Summand}
The investigation begins with the General Zeta Sum Rule, $S(z)=\sum_{k=1}^{\infty}\frac{(z^{k}-1)\zeta(k+1)}{(z+1)^{k}}$. The summand can be decomposed into the difference of two geometric progressions:
\begin{equation}
    \frac{z^k-1}{(z+1)^k} = \left(\frac{z}{z+1}\right)^k - \left(\frac{1}{z+1}\right)^k
\end{equation}
This allows the entire sum $S(z)$ to be expressed as the difference of two Zeta-Geometric series:
\begin{equation}
    S(z) = \Z_1\left(\frac{z}{z+1}\right) - \Z_1\left(\frac{1}{z+1}\right)
\end{equation}
The core task is to understand the properties of the function $\Z_1(x)$.

\section{Application of the Dynamic Recurrence}
The Law of the Golden Recurrence, a consequence of the Cayley-Hamilton theorem applied to the transformation matrix $G$, dictates that the eigenvalues of $G$ (such as $x = \LG = \T+i\J$) must satisfy a linear recurrence. Let $c_1 = \text{Tr}(G) = 2\T$ and $c_2 = -\det(G) = -(\T^2+\J^2)$. The recurrence for the powers of $x$ is:
\begin{equation}
    x^k = c_1 x^{k-1} + c_2 x^{k-2}, \quad \text{for } k \ge 2
\end{equation}
To apply this constraint to our series $\Z_1(x)$, we first separate the $k=1$ term, as the recurrence is not valid for it.
\begin{equation}
    \Z_1(x) = \zeta(2)x + \sum_{k=2}^{\infty} \zeta(k+1)x^k
\end{equation}
We now substitute the recurrence relation (3) into the summation:
\begin{align}
    \Z_1(x) &= \zeta(2)x + \sum_{k=2}^{\infty} \zeta(k+1)\left(c_1 x^{k-1} + c_2 x^{k-2}\right) \\
    &= \zeta(2)x + c_1 \sum_{k=2}^{\infty} \zeta(k+1)x^{k-1} + c_2 \sum_{k=2}^{\infty} \zeta(k+1)x^{k-2}
\end{align}

\section{Re-indexing and Identification of the Emergent Structure}
The two resulting sums must be re-indexed to understand their relation to the original series.
\begin{itemize}
    \item For the first sum, let $j=k-1$. As $k$ runs from $2$ to $\infty$, $j$ runs from $1$ to $\infty$:
    \begin{equation*}
         \sum_{k=2}^{\infty} \zeta(k+1)x^{k-1} = \sum_{j=1}^{\infty} \zeta(j+2)x^{j}
    \end{equation*}
    \item For the second sum, let $m=k-2$. As $k$ runs from $2$ to $\infty$, $m$ runs from $0$ to $\infty$:
    \begin{equation*}
        \sum_{k=2}^{\infty} \zeta(k+1)x^{k-2} = \sum_{m=0}^{\infty} \zeta(m+3)x^{m} = \zeta(3) + \sum_{m=1}^{\infty} \zeta(m+3)x^{m}
    \end{equation*}
\end{itemize}
This reveals that the recurrence does not relate $\Z_1(x)$ to itself, but to a hierarchy of related functions. We define this family of functions as:
\begin{equation}
    \Z_p(x) = \sum_{k=1}^\infty \zeta(k+p)x^k
\end{equation}
In this language, our re-indexed sums become $c_1 \Z_2(x)$ and $c_2\zeta(3) + c_2x\Z_3(x)$.

\section{The Functional Equation}
By substituting the re-indexed forms back into equation (6), we arrive at the final functional equation. This equation is a necessary consequence of imposing the Golden Algebra's dynamical laws onto the Zeta-Geometric series.

\textbf{Theorem.} For a complex number $x$ that is an eigenvalue of the matrix $G$, the family of Zeta-Geometric functions $\Z_p(x)$ obeys the following recurrence relation:
\begin{equation}
    \Z_1(x) = \zeta(2)x + c_2\zeta(3) + c_1 \Z_2(x) + c_2 x\Z_3(x)
\end{equation}
where $c_1 = 2\T$ and $c_2 = -(\T^2+\J^2)$.

This result provides the fundamental structural tool for the advanced analysis of these harmonic Zeta series.

\end{document}